% ============================================================
% Chapter 3: Project Management
%   Section: Budget
%     Subsection: Infrastructure Costs
% Included from chapters/03-project-management/05-budget/budget.tex
% Sources: docs/decisions/deployment.md; backend/data/evaluation/prices.json.
% ============================================================

\subsection{Infrastructure Costs}

\paragraph{Cloud.} The production deployment is sized for one Google
Cloud Compute Engine virtual machine, an \texttt{e2-standard-4} (4 vCPU,
16~GB) in Madrid (\texttt{europe-southwest1}), at about 115~USD a month
on demand or about 73~USD with a one-year commitment, plus a 60~GB
balanced disk and daily snapshots kept seven days
(Section~\ref{sec:dev-deployment}). The size comes from measurements:
about 9.6~GiB were resident across the production stack, and the
language model needs about 45~s per claim on a CPU. The infrastructure
was checked against a local emulator and never applied to a billed
project, so its cost to date is zero; running it for one month of
demonstration would cost about 115~USD.

\paragraph{Language-model providers.} The local model has no per-token
cost: its cost is the machine's time, which the evaluation reports
beside every run (36.5~s per claim at the median in the X-Fact pilot).
Hosted providers are priced per run from a table in the repository
(\path{backend/data/evaluation/prices.json}), filled in from the
provider's pricing page on the day of the run, never from memory. No
hosted model has been run, so this cost is also zero to date.

\paragraph{Search.} The self-hosted metasearch engine has no cost. An
API-keyed search engine, the most effective remedy for refused searches,
is future work precisely because it would add a recurring cost and a
credential (Section~\ref{sec:future-work}).
